\section{Introduction}
\label{sec:introduction}

The rapid expansion of digital payment ecosystems has catalyzed a global surge in electronic transactions, bringing unprecedented convenience to consumers and merchants. However, this growth has concurrently amplified the threat of financial fraud, which inflicts billions of dollars in losses annually on the global economy. In the context of credit card transactions, fraud detection systems serve as the first line of defense for financial institutions. Despite decades of research, effectively identifying fraudulent transactions in real-time remains a critical challenge due to the dynamic and highly skewed nature of the data.

Unlike traditional static datasets, credit card transactions arrive as continuous, high-velocity data streams. Fraud detection models must process and classify these streams sequentially under strict latency constraints. In such dynamic environments, fraud detection faces two paramount challenges. The first challenge is \textit{concept drift}. Fraudsters constantly evolve their tactics to bypass existing security measures, causing the underlying data distribution to shift over time. A static machine learning model deployed in a dynamic environment quickly suffers from performance degradation as its learned patterns become obsolete. To combat this, online learning models equipped with drift detection mechanisms, such as Adaptive Windowing (ADWIN) and Hoeffding Trees, have been proposed to continuously update their knowledge base from incoming data \cite{insert_drift_ref}.

The second, and arguably more insidious challenge, is \textit{extreme class imbalance}. In real-world payment networks, legitimate transactions overwhelmingly dominate the data stream, while fraudulent activities constitute a minuscule fraction (often less than 0.2\%). Traditional online learning algorithms are designed to maximize global accuracy. Consequently, when faced with extreme class imbalance, these models inherently exhibit a strong bias toward the majority class (legitimate transactions), effectively ignoring the minority class (fraud). While techniques like Synthetic Minority Over-sampling Technique (SMOTE) or algorithmic cost-weighting have proven successful in batch learning scenarios, applying them to continuous, high-velocity data streams poses severe computational bottlenecks. 

Existing literature reveals a significant research gap at the intersection of concept drift and extreme class imbalance. Most state-of-the-art streaming algorithms successfully address concept drift but fail to catch fraud when the minority class ratio drops below 1\%. Conversely, methods that attempt to handle both issues typically rely on two-stage preprocessing architectures (e.g., buffering data to oversample before training), which inherently contradict the single-pass requirement of online learning and introduce unacceptable computational latency.

To address these limitations, this paper proposes a novel framework: the Cost-Sensitive Adaptive Random Forest (CS-ARF) for real-time fraud detection. By integrating an \textit{Implicit Online Oversampling} mechanism directly into the objective function of the Hoeffding Trees within the Adaptive Random Forest ensemble, our method seamlessly addresses extreme class imbalance without requiring external preprocessing buffers. The proposed model evaluates incoming transactions by assigning an asymmetric penalty (cost) proportional to the imbalance ratio, ensuring high fraud recall while maintaining the robustness of the ADWIN drift detector. 

The main contributions of this paper are summarized as follows:
\begin{itemize}
    \item We propose a single-pass streaming architecture, CS-ARF, which solves both concept drift and extreme class imbalance simultaneously without relying on computationally expensive two-stage resampling.
    \item We introduce an Implicit Online Oversampling mechanism based on asymmetric cost matrices integrated into the Hoeffding Tree node-splitting criteria.
    \item We present an extensive hyperparameter sensitivity analysis demonstrating the optimal cost boundary for maintaining geometric mean (G-Mean) stability.
    \item We validate our approach on a highly imbalanced credit card fraud data stream, demonstrating a significant increase in fraud recall compared to standard streaming baselines while maintaining low false-positive rates.
\end{itemize}

The remainder of this paper is organized as follows: Section \ref{sec:related_works} reviews related work in online fraud detection. Section \ref{sec:methodology} details the proposed CS-ARF architecture and the Implicit Online Oversampling mechanism. Section \ref{sec:experiments} describes the experimental setup, dataset, and evaluation metrics. Section \ref{sec:results} presents the results and discussion, followed by the conclusion in Section \ref{sec:conclusion}.
